\subsection{The height input and the exact pole predicate}

Put
\[
 S_* =\{p:p\le76176\},\qquad
 \mathcal D=\{0,1,2401,-343/32\},
\]
\[
 H_*=2(4608^{165888}+6^{100})
          \left(\prod_{p\le76176}p\right)^{24}+128.
\]
For $a\in\Q\setminus\mathcal D$, let $M_5(a)$ be the product of
the primes outside $S_*$ at which the point $(a:1)$ has positive odd
contact order with one of $\infty,0,1,2401,-343/32$. Explicitly, choose
local homogeneous coordinates $(u:v)$ and subtract
$\min(v_p(u),v_p(v))$ from the valuations of $v$ and of
$u-\delta v$ for $\delta\in\mathcal D$.

\begin{theorem}[The five-contact height estimate]\label{hyp:o8-height}
For every $a\in\Q\setminus\mathcal D$,
\[
 h(a)\le H_*M_5(a)^{24}.
\]
The displayed constant satisfies $\log_2H_*<4783227$.
\end{theorem}
\begin{proof}
Write $s=a$, $M(s)=M_5(a)$, $H=H_*$,
$B_*=\prod_{p\le76176}p$, and
$\mathcal C=\{\infty,0,1,2401,-343/32\}$.
We prove the geometric comparison, the non-CM estimate through the
modular level, and the CM estimate over the same fiber field.
All heights use the original stable Faltings normalization.
\paragraph{Recovering the contact line from published quotients}
Let $B/\mathbb Q$ be the indefinite quaternion algebra of discriminant $210$,
$\mathcal O$ a maximal order, $X$ its maximal-level Shimura curve, and $W$
its Atkin--Lehner group. The standard genus formula gives
\[
 g(X)=1+\frac{\varphi(210)}{12}=5,
\]
because the elliptic-point terms vanish: $5$ splits in $\mathbb Q(i)$ and
$7$ splits in $\mathbb Q(\sqrt{-3})$, so neither quadratic field embeds in
$B$. The group $W$ is $(\mathbb Z/2\mathbb Z)^4$ with involutions $w_m$
indexed by square classes of divisors of $210$.

The published models of \cite[Theorem 4.2 and Table 2, pp.~938--939]{GonzalezRotger2006} have the following coefficient rows
$(a,b,c,d,e)$ for $y^2=ax^4+bx^3+cx^2+dx+e$:
\[
\begin{array}{c|rrrrr}
 m&a&b&c&d&e\\\hline
30&-43&-686&-2915&-1372&-172\\
42&-43&600&-986&-600&-43\\
70&-43&256&-130&-768&-387\\
105&-43&2&505&12&-1548
\end{array}
\]
The cited theorem identifies each with $X/\langle w_m\rangle$.
The genus formula and the Atkin--Lehner group law used here also appear
in the same published paper, Proposition 5.2 and the paragraph following
Proposition 5.3, pp.~940--941.

Here is a deduction of all the required branch information, rather than
an additional assumption about it. Each of these four genus-one quotients
has eight fixed points upstairs, by Riemann--Hurwitz. A stabilizer for the
$W$ action on a characteristic-zero smooth curve is cyclic, and therefore
has order at most two. Each nonempty fixed set of an involution is a union
of $W$-orbits of size eight. Riemann--Hurwitz for $X\to X/W$ now gives
\[
 8=16(2g(X/W)-2)+\sum_{1\ne w\in W}\#\operatorname{Fix}(w).
\]
The four known fixed sets already contribute $32$. Thus $g(X/W)=0$,
the total is $40$, and there is exactly one further orbit of eight fixed
points. Inertia generators around the branch points have product one.
The four classes $30,42,70,105$ are independent in $W$, and their product
is the class $210$. Hence there are exactly five branch points, each of
index two, with distinct labels $30,42,70,105,210$. Since the involutions
are defined over $\mathbb Q$ and each label occurs only once, each branch
point is rational. Three of them trivialize the genus-zero quotient.

Normalize $b_{210}=\infty$, $b_{105}=0$, $b_{70}=1$, and write
$b_{42}=\lambda$, $b_{30}=\mu$. On a genus-one quotient the residual
group $(\mathbb Z/2)^3$ has translation subgroup of order four, necessarily
the full two-torsion group after choosing an origin. Its quotient has the
same Jacobian $j$-invariant: it is multiplication by two over the algebraic
closure. The remaining double cover of $X/W$ branches at the four points
other than $b_m$. Therefore the four quartic invariants determine the four
Legendre invariants of these four-point subsets.

For a quartic put
\[
 I=12ae-3bd+c^2,\quad
 J_q=72ace+9bcd-27ad^2-27b^2e-2c^3,\quad
 j=\frac{6912I^3}{4I^3-J_q^2}.
\]
The exact computation gives the Legendre parameters
$1/2401,32/375,1/225,5/1029$, in row order. Enumerating the six cross-ratio
values of each of the first two parameters and testing the last two
invariants leaves exactly
\[
 \lambda=2401,\qquad\mu=-343/32.
\]
The script \texttt{check\_height.py} reproduces this rational computation.
Every nonzero determinant of the five primitive contact forms has prime
support in $\{2,3,5,7\}$.

\paragraph{Integral geometry and the degree 576 field}
Choose $\mu_0\in\mathcal O$ with $\mu_0^2=-210$. The embedding criterion
allows this choice in a maximal order; maximal orders are conjugate here.
Use the maximal moduli stack $\mathscr X$ and the entire principal-level
three scheme $Y$, including its geometric components. Yuan's integral
moduli theorem \cite[Theorem 1.1 and Sections 2.1--2.2]{Yuan2024} makes $\mathscr X$ regular, proper and semistable over
$\mathbb Z$, with a universal abelian surface, and smooth away from $210$.
Its generic stabilizer is precisely $\{\pm1\}$. Principal level three is
fine and gives a universal full three-torsion marking.

Work first over $R_0=\mathbb Z[1/210]$. The effective level-change group
has order $36$, prime to every residue characteristic of $R_0$. Its generic
action is free. For a finite tame group acting on a smooth relative curve,
the fixed locus of an element is smooth locally along its fixed points;
an isolated special fixed point consequently lifts. Thus no new fixed
point occurs here, and $Y_{R_0}\to X_{R_0}$ is finite etale. The $W$ action
extends to the smooth stable model of $X$. Its stabilizers remain cyclic,
of order at most two. Distinct involution fixed loci cannot intersect;
each fixed locus is etale over the base. Consequently $X_{R_0}/W$ is a
smooth genus-zero curve with five disjoint branch sections of index two.
The three rational sections already chosen identify it with
$\mathbb P^1_{R_0}$. Their values are the five elements
of $\mathcal C$.

Set $K=\widehat{\mathcal O}^{\times}$, $U=U(3)$,
$Z=\mathbb A_f^\times$, and $N=N_{B(\mathbb A_f)^\times}
(\widehat{\mathcal O})$. Normalization preserves $3\widehat{\mathcal O}$,
so $U$ is normal in $N$. The canonical adelic action is defined over the
reflex field $\mathbb Q$ \cite[2.2.5 and 2.3.1]{Deligne1979}; this refers to the entire canonical model, not
an individually selected complex component. Every central finite idele
is a rational scalar times a unit idele. At three its scalar residue is
$1$ or $-1$, which a rational sign removes. Thus $Z\subset\mathbb Q^\times U$
acts trivially on $Y$.

At three the quaternion algebra is ramified. Its maximal order modulo
three has $81$ elements, radical of size $9$, and hence $72$ units.
It follows that
\[
 1\longrightarrow (\mathcal O/3\mathcal O)^\times/\{\pm1\}
 \longrightarrow\Gamma=N/(ZU)\longrightarrow W\longrightarrow1,
 \qquad |\Gamma|=36\cdot16=576.
\]
Above an unbranched point the actions of $W$ and of the effective
level-change kernel are free. The fiber therefore is a torsor of size
$576$ for this constant group. The Galois permutations commute with its
regular action, so belong to its regular centralizer, also of order
$576$. The normal splitting field $L$ consequently satisfies
$[L:\mathbb Q]\mid576$: the Galois group is a subgroup of that
regular centralizer. In particular $[L:\mathbb Q]\le576$.

All fiber surfaces and their three-torsion markings are defined over $L$.
For two such surfaces, a homomorphism annihilating the three-torsion
factors through $[3]$. Hence reduction embeds
\[
 \operatorname{Hom}(A,B)/3\operatorname{Hom}(A,B)
 \hookrightarrow\operatorname{Hom}(A[3],B[3]).
\]
The finite Galois image on the integral Hom lattice is trivial modulo
three. The principal congruence subgroup modulo three in an integral
general linear group is torsion-free, so this image is trivial. Thus
every geometric homomorphism between fiber surfaces is defined over $L$.

At a prime outside $210$, tame local coordinates for the quotient are
$w=z^2$ up to a unit at a branch section, and etale elsewhere. A base
point with even contact gives no ramification after an unramified
extension. The universal abelian scheme extends at the resulting
integral point of $Y_{R_0}$. Therefore
\[
 \operatorname{rad}(D_LN_{\rm bad})\mid B_*M(s).
\tag{R}\label{eq:h5-R}
\]
This assertion includes contact with infinity. It is not a statement
obtained merely from the numerical prime support of pairwise determinants.

The normalizer correspondences give isogenies between the conjugate
surfaces. In the non-CM case the geometric endomorphism algebra is $B$;
Skolem--Noether adjusts a conjugate isogeny to commute with its chosen
quadratic subfield. All such maps lie over $L$ by the Hom argument.
These facts establish the usual GL$_2$-type and compatibility hypotheses;
they do \emph{not} establish the separate bounded descent to a
GL$_2$-type variety over $\mathbb Q$.

\paragraph{The metric comparison, with its exact additive budget}
Put $t=32s$, $(a_1,a_2,a_3,a_4)=(0,32,76832,-343)$ and
$P(t)=\prod_i(t-a_i)$. The sections
\[
 q_0=(dt)^2/P(t),\qquad q_1=tq_0
\]
are holomorphic bicanonical sections on $X$, with no common zero. At a
finite branch point use $t-a_i=z^2$; at infinity use $t=z^{-2}$.
At infinity $q_0$ vanishes to order two and $q_1$ does not vanish.

For each prime let $\delta_p=\max_{i<j}v_p(a_i-a_j)$, so
\[
 J=\prod_p p^{\delta_p}=2^{10}3^2 5^3 7^4=2765952000.
\]
Both $J^2q_i$ extend across the maximal regular stack. Indeed, at the
generic point of a vertical component semistability makes $p$ a
uniformizer, and the relative differential satisfies $v(du)\ge v(u)$.
Writing $r_i=v(t-a_i)$ gives $v(dt)\ge\max r_i$. If $v(t)<0$, all $r_i$
equal $v(t)$ and both sections have nonnegative valuation. Otherwise all
$r_i\ge0$ and at most one exceeds $\delta_p$. In either case
$2\max r_i-\sum r_i\ge-2\delta_p$. This proves the required vertical
integrality; regularity handles codimension two.

For the archimedean estimate write $r_i=|t-a_i|$, $U_0=1+\sum_i r_i$ and
\[
 ds_\rho^2=\frac{U_0}{\prod_i r_i}|dt|^2,\qquad
 C_{\rm met}=16(12\cdot76832+9)^3=12540153541944538512.
\]
Its pullback is smooth and positive at every branch point and infinity.
The curvature formula, $\Delta U_0=\sum_i r_i^{-1}$ and
$|\nabla U_0|^2\le16$ give
\[
 -K_\rho\ge\frac{\sum_i\prod_{j\ne i}r_j}{2U_0^3}.
\]
When $|t|\le2\cdot76832+2$, at most one $r_i<1/2$, and
$U_0\le12\cdot76832+9$. The displayed bound is at least
$C_{\rm met}^{-1}$. Outside this disk, $|t|/2\le r_i\le3|t|/2$
and $U_0\le7|t|$, giving the stronger bound $1/1372$.
The maximum principle against the curvature $-1$ metric implies
$ds_\rho^2/C_{\rm met}\le ds_{\rm hyp}^2$. Yuan's Petersson norm is
$\|d\tau\|=2\operatorname{Im}\tau$. Consequently
\[
 \|q_0\|_{\rm Pet}\le4C_{\rm met}/U_0,\qquad
 \|q_1\|_{\rm Pet}\le4C_{\rm met}|t|/U_0\le4C_{\rm met}.
\]

Yuan's isometry \cite[Theorem 1.1 and metric definitions, pp.~2--3]{Yuan2024} identifies the square of the Hodge bundle with
$210\omega_{\mathscr X}^{\otimes2}$ with the induced metric. The Hodge
height in his normalization equals the original stable Faltings height
plus $\log\pi$ in dimension two. Thus
\[
 h_{\bar\omega^{\otimes2}}(y)=2h_F(A_y)+2\log\pi+\log210.
\]
The sign of $\log210$ is positive: the sublattice $210\omega^{\otimes2}$
has height $h_{\bar\omega^{\otimes2}}-\log210$. Evaluate on an extension
where the point extends to the proper maximal stack; its universal
abelian scheme computes stable height. The integral sections and the
supremum estimate yield
\[
 h(s)\le2h_F(A_y)+2\log\pi+\log210+2\log J
                       +\log(4C_{\rm met})+\log32.
\]
The integer comparison
$210J^2(4C_{\rm met})32\cdot16<2^{160}$, together with $\pi<4$ and
$\log2<3/4$, bounds the additive term by $120<128$. Therefore
\[
 h(s)\le2\max(0,h_F(A_y))+128.
\tag{G}\label{eq:h5-G}
\]


% Proposed repair of the original H5 arithmetic step. Independent review pending.
% Uses the existing baseline bibliography keys, plus GaudronRemond2023 and
% the published 2021 edition of Voight, chapter 43.
\paragraph{A level bound for the non-CM fibers}
\label{sec:h5-level-repair}

This argument allows the scalar character order and coefficient field to
vary. It does not assert bounded same-field descent to a rational
GL$_2$-type variety. Let $s$ be a rational point outside the five contacts,
let $L/\mathbb Q$ be the normal fiber field constructed above, and let
$A/L$ be a non-CM fiber surface. Thus $[L:\mathbb Q]\mid576$, all
geometric homomorphisms between conjugate fibers are defined over $L$,
and $\operatorname{End}^0(A_{\overline{\mathbb Q}})=B$, the quaternion
division algebra of discriminant $210$. Put
\[
 R=B_*M(s),\qquad T=\{p:p\mid R\}.
\]
The geometric argument gives that $L$ is unramified outside $T$ and
$A$ has good reduction outside the primes above $210$. All conjugate
QM structures lie in the same normalizer orbit.

\paragraph{Good reduction after full torsion.}
First a QM surface with the indicated division action has potentially
good reduction at every finite place. Indeed, pass to an extension with
semistable reduction, using \cite[Theorem 3.4, p.~406]{silverbergzarhin}.
Every endomorphism extends uniquely to the N\'eron model by its mapping
property. On the identity component of the geometric special fiber it
preserves the maximal torus $\mathcal T$: this torus is the maximal
connected affine subgroup, and every homomorphism from a torus to an
abelian variety is zero. If its rank $t$ is positive, restriction gives
a unital homomorphism
\[
 B\longrightarrow\operatorname{End}^0(\mathcal T)=M_t(\mathbb Q),
 \qquad 1\le t\le2.
\]
The kernel is a two-sided ideal of the simple algebra $B$, and the map
is unital, so it is injective. For $t=1$ dimensions prohibit this;
for $t=2$ it would be an isomorphism from a division algebra to
$M_2(\mathbb Q)$, also impossible. Hence $t=0$. The semistable identity
component is abelian, which is good reduction. These N\'eron-model and
reduction conventions are those of \cite[\S3, pp.~405--406]{silverbergzarhin}.

Set $E=L(A[11])$. We next justify that $E/\mathbb Q$ is Galois.
For each $\sigma$, the normalizer relation supplies a compatible
isogeny ${}^{\sigma}A\to A$ of degree prime to $11$. Here is the
prime-to-$11$ detail. In the analytic description
$A_z=\mathbb C^2/(\mathcal O(z,1)^t)$ of the published 2021 edition
of Voight \cite[equations (43.1.2)--(43.1.3)]{voight2021}, a positive normalizer element
$b=\left(\begin{smallmatrix}a&b'\\c&d\end{smallmatrix}\right)$
gives the scalar compatible map $(cz+d)I:A_{bz}\to A_z$.
When $\mathcal O b\subset\mathcal O$, its degree is
$[\mathcal O:\mathcal O b]=\operatorname{Nrd}(b)^2$.
At $11$ the algebra is split and the maximal-order normalizer is
$\mathbb Q_{11}^*\operatorname{GL}_2(\mathbb Z_{11})$.
Rescale $b$ rationally to make it an $11$-adic unit, then clear the
remaining denominators by an integer prime to $11$. This produces
the claimed isogeny. Changes of level marking use isomorphisms.
All these geometric homomorphisms are $L$-defined by construction.
Consequently the $11$-torsion representations of $A$ and its conjugates
over $L$ have the same kernel. Thus every conjugate of $E$ equals $E$.

The faithful action on $A[11]$ now gives
\[
 [E:\mathbb Q]\le576|\operatorname{GL}_4(\mathbb F_{11})|
 <2^{10}11^{16}<2^{66}.
\tag{LR1}
\]
For residue characteristic different from $11$, full rational
$11$-torsion gives semistable reduction over $E$ by
\cite[Theorem 3.5, p.~406]{silverbergzarhin}. The preceding torus
argument makes this good reduction. In particular it applies at
each of $2,3,5,7$.

\paragraph{The scalar lift with controlled support.}
Choose $B$-compatible isogenies $\mu_\sigma:{}^\sigma A\to A$ over $L$,
with $\mu_1=1$, and inflate this finite choice to $G_{\mathbb Q}$.
Their discrepancies $c(\sigma,\tau)$ commute with $B$, so lie in
$\mathbb Q^*$. Put $a_\sigma=(\deg\mu_\sigma)^{-1/4}>0$ and
$\nu_\sigma=a_\sigma\mu_\sigma$, with trivial Galois action on the
added algebraic coefficients. Since a scalar $c$ has degree $c^4$,
the discrepancies of the $\nu_\sigma$ form a sign cocycle
$\varepsilon\in Z^2(G_{\mathbb Q},\{\pm1\})$, inflated from
$\operatorname{Gal}(L/\mathbb Q)$.

The sign class is unramified, hence zero, locally outside $T$.
At each odd prime where its Brauer invariant is $1/2$, choose an odd
Dirichlet character modulo $p$; at $2$ use the odd character modulo
$4$ if required. Their product $\chi$ is tame at every odd prime,
has local conductor at most $4$ at $2$, and has order
\[
 m\le2\prod_{\substack{p\in T\\p\ne2}}(p-1)\le2R.
\]
For constant coefficient groups, the boundary in
$1\to\{\pm1\}\to\mu_\infty\xrightarrow{2}\mu_\infty\to1$
is the obstruction to a character square root. Local reciprocity
identifies this obstruction with the value at $-1$: the cyclic
torsion factor has precisely this obstruction, while the valuation
and principal-unit factors admit character square roots.
Thus the boundary of $\chi$ equals $[\varepsilon]$ at every finite
place. Brauer reciprocity and the global-to-local injection give
equality globally. A pointwise square root, corrected by a sign
cochain, gives
\[
 b:G_{\mathbb Q}\longrightarrow\mu_{2m},\qquad
 b^2=\chi,\qquad\delta b=\varepsilon.
\tag{LR2}
\]
Only this \emph{sign} cochain is asserted to take root-of-unity values;
the full scalar splitting also involves the $a_\sigma$.

On $G_L$, $b$ is a conjugation-invariant character. Outside $T$ its
inertia characters are quadratic. Only finitely many are nontrivial;
twist by the quadratic character of the product of those odd primes
to kill them. At each odd prime there is a unique nontrivial quadratic
inertia character. This changes no square or coboundary in (LR2),
and introduces possible ramification only at $2\in T$.
The kernel field $L'$ of $b|G_L$ is normal over $\mathbb Q$ and
\[
 [L':\mathbb Q]\le2m[L:\mathbb Q]\le2304R.
\tag{LR3}
\]

\paragraph{The compatible modular representations.}
Choose a number field containing all normalization scalars, the values
of $b$, and a splitting field of $B$. For each coefficient prime $\ell$
and an embedding into $\overline{\mathbb Q}_\ell$, set
\[
 \rho_\ell(\sigma)=b(\sigma)^{-1}a_\sigma\mu_\sigma\circ\sigma.
\]
The cocycle identity makes this a representation commuting with $B$;
after splitting $B$, matrix units give $\rho_\ell=r_\ell\oplus r_\ell$
with $\dim r_\ell=2$. On $G_{L'}$, $\rho_\ell$ is the Tate action of
$A$. Faltings' semisimplicity and endomorphism theorems and $\operatorname{End}^0(A)=B$
show that the multiplicity representation is absolutely irreducible.
Proper-smooth comparison and finite descent, as in
Lemma~\ref{lem:de-rham-descent}, make $r_\ell$ de Rham with Tate
Hodge--Tate weights $\{0,-1\}$; Berger's monodromy theorem makes it
potentially semistable.

It is odd. On Betti homology complex conjugation followed by the
holomorphic isogeny $\mu_c$ exchanges the two Hodge summands, each
a simple $B\otimes\mathbb C$-module. Its action on the multiplicity
space is therefore nonscalar. After the scalar correction its square
is one, so its eigenvalues are $1,-1$. The representation is unramified
outside $T\cup\{\ell\}$, by (LR2) and the reduction and field properties.
Pan's regular two-dimensional modularity theorem applies at $\ell=5$.
We use its published formulation in \cite[Section 7, p.~624]{calegari-icm2022},
with the definition of a geometric representation on pp.~611--612,
and the detailed author proof in \cite[Theorem 1.0.4]{pan2022}.
It requires no additional residual-image restriction at this prime and
identifies $r_5$ with a cuspidal modular representation up to a character
twist. The published formulation is used as stated; no byte or text
identity between Pan's author version and the journal version is asserted.

The twist can be taken finite and the weight is two, without invoking
a separate local-algebraicity theorem for characters. The twisting
character is a de Rham subrepresentation of the relevant Hom
representation. Comparing its integer weight with $\{0,-1\}$ and
the modular weights $\{0,1-k\}$ forces $k=2$ and twist weight zero.
On $G_{L'}$, the identity
$(\det r_5)^2=\det V_5(A)=\chi_{\rm cyc}^2$ shows that
$\det r_5/\chi_{\rm cyc}$ has finite image, since $G_{L'}$ has finite
index in $G_{\mathbb Q}$. The modular determinant
divided by $\chi_{\rm cyc}$ also has finite image. The square of the
twisting character consequently has finite image, so the character
itself has finite image. Absorb it in the modular form and take the
primitive weight-two newform $f$ of level $N$.

The conductor at $5$ is not omitted. For each prime
$p\notin T\cup\{3,5\}$, reduction of $\mu_\sigma$ composed with
relative $p$-Frobenius is a quasi-endomorphism of the good reduction
of $A$ at a chosen place. After clearing denominators its characteristic
polynomial is independent of the Tate prime. Half its trace, multiplied
by $b(\sigma)^{-1}a_\sigma$, is therefore one algebraic number giving
the trace of every $r_\ell$. Enlarge the coefficient field to contain
the coefficients of $f$. The equality with the trace of $f$ at five
is an equality of algebraic numbers; transporting it to an embedding
at three and applying Chebotarev and semisimplicity identifies $r_3$
with the three-adic realization of $f$. Thus use $r_3$ to compute the
conductor at $5$ and $r_5$ at every other prime. The modular level is
the conductor of these realizations by
\cite[\S\S0.3--0.7, Theorem A, pp.~409--411]{carayol}
and newform theory \cite[Theorems 1 and 4]{casselman}.
Here the totally real field is $\mathbb Q$, of odd degree one, so
Carayol's extra finite-place hypothesis for even degree does not apply.
His unramified normalization twist and contragredient preserve the
conductor; his theorem allows arbitrary nebentypus.
Li's Fricke relation and Shimura's
Jacobian construction, as in the baseline modularity argument, also
place $r_5$ in the Tate module of $J_1(N)$.

\paragraph{The primitive level equals the local conductor.}
Let $f\in S_2(\Gamma_0(N),\chi)$ be the primitive newform obtained
after absorbing the finite character twist, and let $\Pi_f$ be its
unitary cuspidal automorphic representation. The character $\chi$
is arbitrary; its conductor need only divide $N$.
For every prime $p$, write $a(\Pi_{f,p})$ for the local conductor
exponent. Then
\[
 a(\Pi_{f,p})=v_p(N).
\tag{CL1}
\]
Here is the exact published newform input. In
\cite[Theorem~1, p.~302]{casselman}, the local conductor ideal is
the largest ideal $\mathfrak c$ for which the appropriate central
character eigenspace under $\Gamma_0(\mathfrak c)$ is nonzero;
at that ideal the eigenspace is one-dimensional. The remark on
pp.~306--307 identifies this ideal with the exponent in the local
epsilon factor. Theorem~4 on p.~313 gives the global level spaces:
they vanish unless $\mathfrak c(\Pi)\mid M$, and at
$M=\mathfrak c(\Pi)\prod_p p^{b_p}$ their dimension is
$\prod_p(b_p+1)$. The paragraph immediately following that theorem
identifies all primitive newforms with the minimal level
$M=\mathfrak c(\Pi)$. This proves (CL1).

This statement includes nebentypus. Casselman's classical setup on
p.~309 allows every character $\epsilon$ modulo $M$ and uses
$f|_k\gamma=\epsilon(a)^{-1}f$ for
$\gamma=\left(\begin{smallmatrix}a&b\\c&d\end{smallmatrix}\right)
\in\Gamma_0(M)$. Since $ad\equiv1\pmod M$, this is exactly
$\epsilon(d)f$. Thus the usual nebentypus convention is already
covered. No squarefree-level hypothesis occurs.

For a coefficient characteristic $\ell\ne p$, Carayol's published
Theorem~A gives the full local Weil--Deligne compatibility, not only
an equality at unramified primes
\cite[\S\S0.3--0.7, pp.~409--411]{carayol}.
Its base field here is $\mathbb Q$, its real weight is two, and its
extra finite discrete-series-place condition applies only to even
degree totally real fields, so imposes no additional hypothesis
here. Theorem~A includes the dyadic cases; the weaker Theorem~B
is not the result being used.

Carayol normalizes reciprocity by sending geometric Frobenius to
a uniformizer. His Hecke correspondence in \S0.5 differs from
local Langlands by a contragredient and an unramified norm twist.
It is characterized there by equality of local $L$- and
epsilon-factors, including character twists. The conductor is
unchanged by either duality or an unramified twist. In particular,
the passage to the Tate-module convention used for $r_\ell$ changes
no conductor exponent. Equality of good Frobenius traces identifies
the global representation supplied by Carayol with the already
constructed irreducible modular representation, in this convention.
This uses Chebotarev and Brauer--Nesbitt; it does not use the
unpublished irreducibility observation in Carayol's remark after
Theorem~A. Since the matching representation is irreducible, its
matching semisimplification also forces the original representation
to be irreducible.

The epsilon-factor conductor in (CL1) therefore equals the Artin
conductor of $r_\ell|_{G_{\mathbb Q_p}}$:
\[
 a_p(r_\ell)=a(\Pi_{f,p})=v_p(N),\qquad p\ne\ell.
\tag{CL2}
\]
The Artin conductor here includes its tame term and Swan term;
equivalently one uses the full Weil--Deligne conductor, including
the monodromy contribution. One must not discard monodromy by
replacing the local representation with its semisimplification on
inertia. Carayol's compatibility supplies precisely the required
full local object. His epsilon-factor normalization and
Casselman's conductor-zero additive character give the same
exponent; changing an additive character only introduces the usual
common dimension-dependent shift, so no such change is needed.

Apply (CL2) with $\ell=5$ at every $p\ne5$ and with $\ell=3$ at
$p=5$. The companion proof identifies both $r_5$ and $r_3$ with
the realizations of this same $f$. Thus the asserted conductor
bounds control every prime divisor of the primitive level, even
when $p^2\mid N$. No compatibility theorem at $p=\ell$ is needed.

\paragraph{The conductor bound.}
Let $F=E\mathbb Q(\chi)$, where $\mathbb Q(\chi)$ means the kernel
field of the character, not its value field. On $G_F$, $b$ is a
conjugation-invariant quadratic character. Its kernel gives a Galois
extension $F'/\mathbb Q$ of degree at most two over $F$.
For $p=2,3,5,7$, let $e_E,e_{F'}$ denote local ramification indices
over $\mathbb Q_p$. From (LR1), $v_p(e_E)\le65$.
At odd $p$ both added extensions are tame, whereas at $2$ each
adds ramification index at most two. Consequently
\[
 v_p(e_{F'})\le65\ (p=3,5,7),\qquad v_2(e_{F'})\le67.
\tag{LR4}
\]
The possibly unbounded tame ramification index causes no problem.
For a finite Galois extension of $\mathbb Q_p$ with ramification
index $e$ and different exponent $\delta$, an Eisenstein polynomial
over the maximal unramified subfield gives
\[
 \delta\le e-1+ev_p(e).
\]
Indeed the derivative terms have distinct valuations modulo $e$,
and its leading derivative term has the stated valuation. In upper
numbering the different formula is
\[
 \delta/e=1-1/e+\int_0^\infty(1-|G^u|^{-1})\,du.
\]
Before the last positive break the integrand is at least $1/2$.
Thus (LR4) bounds this break for $F'/\mathbb Q_p$ by $136$.
These ramification and conductor conventions are those of
\cite[\S\S4.1--4.3, pp.~155--159]{serrelocal}.
Over $F'$ the scalar correction is trivial and $A$ has good reduction.
Above that break, therefore, the two-dimensional representation is
trivial on upper wild groups. Its Swan conductor is at most $272$,
and its tame conductor is at most two. Hence
\[
 a_p(r)\le274<300\qquad(p=2,3,5,7).
\]
For $p>7$, $L/\mathbb Q$ is tame because its degree divides $576$,
$A/L$ has good reduction, and $b$ is trivial on wild inertia:
on that pro-$p$ group its square is tame and its remaining image has
order at most two. Thus $a_p(r)\le2$ there; it is zero outside $T$.
Using the coefficient-prime comparison above proves
\[
 N\le210^{300}R^2.
\tag{LR5}
\]

\paragraph{The original constant is preserved.}
Apply the accompanying modular-Jacobian transfer calculation with
$C=210^{300}$ and (LR3)--(LR5). It gives
\[
 h_F(A)<2^{24}210^{3600}R^{24}
 <2^{28824}R^{24}<4608^{165888}R^{24}.
\tag{LR6}
\]
The last comparison uses $210<2^8$, $4608>2^{12}$, and
$28824<1990656=12\cdot165888$. Together with the separately proved
CM bound and geometric height comparison, (LR6) yields exactly
\[
 h(s)\le
 \bigl(2(4608^{165888}+6^{100})B_*^{24}+128\bigr)M(s)^{24}.
\]
No contact, coefficient constant, or radical exponent has been changed.
% UNAPPLIED CANDIDATE: global review and published Pan match are pending.

% Actual level-based repair: no uniform coefficient-field bound is used.
% The modular representation and level estimate are proved in the companion
% height-level-repair fragment. This file proves its height consequence.
\paragraph{Transfer from the level, allowing unbounded coefficient degree.}
Let $A/L$ be the non-CM quaternionic surface. Write $R\geq1$ for the
squarefree product of the prescribed rational bad primes, including $2$,
and assume $[L:\mathbb Q]\leq576$. The modularity construction provides
a weight-two newform of level $N$ and a field $L'$ over which the Tate
modules of $A$ and $J=J_1(N)$ share an irreducible constituent. In the
notation of the sign-splitting construction, the auxiliary Dirichlet
character has order
\[
 m\leq 2\prod_{\substack{p\mid R\\p\ne2}}(p-1)\leq2R.
\]
The sign-splitting cochain has values in $\mu_{2m}$. If one instead
uses the full scalar splitting map, it also includes the square root
of the degree map; that factor equals one on $G_L$. Thus in either
notation the restriction to $G_L$ is a character valued in
$\mu_{2m}$, and its kernel field satisfies
\[
 d'=[L':\mathbb Q]\leq2m[L:\mathbb Q]\leq2304R.
\tag{LT1}
\]
The quadratic twist used to remove extra ramification does not increase
this bound, since $2$ divides $2m$. There is no bound here on $m$
independent of $R$.

Faltings' Hom theorem, followed by scalar extension in the Hom space,
gives a nonzero $L'$-homomorphism $A\to J$. Absolute simplicity of $A$
and Poincar\'e reducibility therefore give an $L'$-isogeny
$A\times Z\to J$ for an abelian variety $Z/L'$. Put $g=\dim J$;
then $g\geq2$. The modular genus estimate
\cite[(5.9), p.~13]{vonKaenel2021} gives $g\leq N^2/24$.
The standard Belyi map has degree at most $N^2$ and
\cite[Theorem~1.1.1 and \S2.3]{javanpeykar} gives
\[
 h_F(J)\leq13000000N^{12}.
\tag{LT2}
\]
All heights in this argument use the original stable Faltings metric.

Suppose the companion local argument has proved $N\leq C R^2$,
where $C\geq1$. Then
\[
 g\leq C^2R^4/24,\qquad
 h_F(J)\leq13000000C^{12}R^{24}.
\]
Set $l=\log C$ and $r=\log R$. The published
Gaudron--R\'emond Theorem~1.9(1) \cite[pp.~7--8]{GaudronRemond2023}, combined with half-logarithmic isogeny-height variation,
gives
\[
 h_F(A)+h_F(Z)\leq h_F(J)+V,
 \quad
 V=2g^2\left(8g^2\log(7g)+\log d'
       +\log\max(1,\log d',h_F(J))\right).
\tag{LT3}
\]
Here the source variety for the isogeny estimate is $J_{L'}$.
The three logarithms in (LT3) satisfy respectively
\[
 \log(7g)\leq2l+4r,\quad
 \log d'<8+r,\quad
 \log\max(1,\log d',h_F(J))<17+12l+24r.
\]
For the last assertion, the bound in (LT2) also dominates $1$ and
$\log d'$; use $8+r\leq9R^{24}$.
It follows that
\[
 V\leq
 \frac{C^8R^{16}}{24^4}(32l+64r)
 +\frac{C^4R^8}{288}(25+12l+25r).
\tag{LT4}
\]
Using $l\leq C^4/4$, $r\leq R^8/8$, and also
$r\leq R^{16}/16$ in the second summand gives
\[
 \frac{V}{C^{12}R^{24}}
 \leq\frac{16}{24^4}+\frac{25+3+25/16}{288}<1.
\tag{LT5}
\]
The Bost lower bound in \cite[(2.2), p.~5]{vonKaenel2021}
gives $h_F(Z)\geq-\tfrac12\dim Z\log(2\pi^2)>-\tfrac32g$.
Consequently
\[
 h_F(A)<(13000000+1+1)C^{12}R^{24}
       <2^{24}C^{12}R^{24}.
\tag{LT6}
\]
This calculation needs no uniform bound on the degree of the newform's
coefficient field and no assertion that all its geometric factors are
isogenous to $A$.

For the companion conductor estimate one may take $C=210^{300}$.
Since $210<2^8$ and $4608>2^{12}$,
\[
 2^{24}C^{12}=2^{24}210^{3600}
 <2^{28824}<2^{1990656}<4608^{165888}.
\tag{LT7}
\]
Thus the non-CM height inequality is strictly within the original
coefficient budget, once the stated modularity and level estimate have
been established. In the five-contact application $R=B_*M(s)$, so
the radical exponent is exactly $24$, as required.

% Original 15671 architecture: independent CM branch, version 1.
% This lemma is conditional on the separately audited field/ramification lemma.
\paragraph{The CM height estimate over the same field}
Let $A/L$ be an abelian surface of CM type, let $B$ be a quaternion
division algebra over $\mathbb Q$ embedded in
$\operatorname{End}^0_{\overline L}(A)$, and suppose that all geometric
endomorphisms of $A$ are defined over $L$. Assume
\[
 d=[L:\mathbb Q]\le576,\qquad
 \operatorname{rad}(|D_L|)\le R,\qquad R\ge1.
\]
All heights below are stable Faltings heights in Faltings's original
normalization. We prove
\begin{equation}\label{eq:original-cm-bound}
 h_{\rm F}(A)<854R<6^{100}R^{24}.
\end{equation}
No assertion about the existence or degree of $L$ is included in this lemma.

First, $A_{\overline L}$ is isogenous to $E^2$ for a CM elliptic curve
$E$. Here is a dimension argument which makes the alternative cases
explicit. If $A$ were geometrically simple, its division endomorphism
algebra would act faithfully on the four-dimensional rational homology.
Its dimension therefore divides four. As it contains $B$, it would equal
$B$; but $B$ has no commutative semisimple subalgebra of dimension four,
contrary to CM type. If the two elliptic factors were not isogenous,
their endomorphism algebra would be a product of commutative fields and
could not contain $B$. Finally, if $E$ had no CM, the endomorphism algebra
would be $M_2(\mathbb Q)$, and the embedding of the four-dimensional
division algebra $B$ would be an impossible isomorphism.

The elliptic factors and the isogeny may be taken over $L$. Indeed,
the rational idempotents in $\operatorname{End}^0(A)$ are defined over
$L$, and the image of an integral multiple of such an idempotent is an
elliptic subvariety over $L$. Endomorphisms and homomorphisms of the
resulting factors extend to rational endomorphisms of $A$ by the
inclusions and rational projections; these are fixed by
$\operatorname{Gal}(\overline L/L)$. Thus the factor isogenies, and all
endomorphisms of $E$, are defined over $L$.

Put $F=\operatorname{End}^0(E)$, an imaginary quadratic field, and
$\mathcal O=\operatorname{End}(E)$. Differentiation gives an injective
unital map $F\longrightarrow\operatorname{End}_L(\operatorname{Lie}E)=L$:
it is nonzero, and every nonzero elliptic endomorphism in characteristic
zero is separable. Consequently $F\subset L$.

We next pass to the maximal order without extending $L$. Choose
$f\ge1$ with $f\mathcal O_F\subset\mathcal O$, choose an integral
basis $\alpha_1,\alpha_2$ of $\mathcal O_F$, and form the finite
subgroup scheme
\[
 K=\sum_{i=1}^2(f\alpha_i)(E[f])\subset E.
\]
Each $f\alpha_i$ is defined over $L$, so $K$ and the quotient
$E'=E/K$ are defined over $L$. Under a complex uniformization
$E(\mathbb C)=\mathbb C/\Lambda$, this subgroup is
$\mathcal O_F\Lambda/\Lambda$; thus $E'(\mathbb C)$ has lattice
$\mathcal O_F\Lambda$ and endomorphism ring $\mathcal O_F$.
The ring cannot be larger, since the rational endomorphism algebra
remains $F$. This is also a direct special case of the same-field
maximal-order theorem of R\'emond, but the construction just given is
the proof used here. It follows that $A'=(E')^2$ is $L$-isogenous to $A$.

Write $\Delta_F$ for the fundamental discriminant of $F$. The
discriminant tower formula implies that every rational prime dividing
$\Delta_F$ divides $D_L$. A quadratic fundamental discriminant has
odd prime exponents one and exponent at two at most three. Therefore
\[
 |\Delta_F|\le4\operatorname{rad}(|D_L|)\le4R.
\]
The multiplier four includes the possible factor $2^3/2$.

For completeness we give an elementary analytic bound on $h_{\rm F}(E')$.
CM elliptic curves have potentially good reduction. After a finite
extension with everywhere good reduction, the elliptic discriminant
formula of Silverman, reproduced as equation (6.2) in
\cite{cmheight}, gives
\[
 h_{\rm F}(E')=\frac1{[L':\mathbb Q]}\sum_{\sigma:L'\hookrightarrow\mathbb C}
 \left(-\log(2\pi)-\log|\eta(\tau_\sigma)|^2
                     -\tfrac12\log\operatorname{Im}\tau_\sigma\right).
\]
Here $L'$ is used only to evaluate the stable height; it is not the
field used in the isogeny estimate below. The discriminant term is zero
because reduction is good. We use the individual-embedding formula
(6.2), not the class-average assertion of Proposition 6.1, so no
hypothesis on Weierstrass coefficients lying in $\mathbb Q(j)$ is needed.

For each embedding choose a reduced quadratic parameter of discriminant
$\Delta_F$. If $y=\operatorname{Im}\tau$, then
\[
 \frac{\sqrt3}{2}\le y=\frac{\sqrt{|\Delta_F|}}{2a}
                   \le\frac{\sqrt{|\Delta_F|}}2,
 \qquad a\in\mathbb Z_{\ge1}.
\]
Put $r=|e^{2\pi i\tau}|<1/16$. The product defining $\eta$ gives
\[
 -\log|\eta(\tau)|^2
 \le\frac{\pi y}{6}+2\sum_{n\ge1}\frac{r^n}{1-r^n}
 \le\frac{\pi y}{6}+\frac{2r}{(1-r)^2}
 <\frac{\pi\sqrt{|\Delta_F|}}{12}+\frac{32}{225}.
\]
Also $-\tfrac12\log y\le\tfrac14\log(4/3)<1/12$.
Discarding the negative term $-\log(2\pi)$ and using $\pi<4$ shows
\[
 h_{\rm F}(E')<\sqrt{|\Delta_F|}+1\le2\sqrt R+1\le3R.
\]
The height is additive on products, hence $h_{\rm F}(A')<6R$.

Apply \cite[Theorem 1.9(1)]{GaudronRemond2023} over $L$, with source
$A'$ and target $A$ and with no extension of the base field.
For a surface, its quantity $\Xi(A')$ is
\[
 \Xi(A')=\left(14^{32}d\max(1,\log d,h_{\rm F}(A'))\right)^8.
\]
The general assertion supplies an isogeny $A'\to A$ of degree at most
$\Xi(A')^2$; no identification of CM type with that paper's
separate ``MM'' condition is made or needed. The isogeny-height
inequality then gives
\[
 h_{\rm F}(A)\le6R+
 8\log\left(14^{32}\,576\max(1,\log576,6R)\right).
\]
One can see the last inequality directly from invariant differentials:
after semistable base change, pullback on determinant differentials has
integral cokernel, while its archimedean squared norm is multiplied by
the isogeny degree. Thus the height difference is half the logarithm
of that degree minus a nonnegative normalized cokernel length.

Finally, $\log14<3$, $\log576<7$, and
$\max(1,\log576,6R)\le7R$. Since $\log7<2$ and $\log R\le R$,
\[
 h_{\rm F}(A)<6R+8(96+7+2+\log R)
             \le854R.
\]
As $6^4=1296>854$ and $R\ge1$, this proves
\eqref{eq:original-cm-bound}. In particular the original
$6^{100}R^{24}$ allowance absorbs this CM estimate with the field
degree still exactly bounded by 576.

\paragraph{The unconditional numerical certificate and scalar detector}
The specified integer $H$ satisfies
\[
 \log_2H<4783227.
\]
Indeed the exact sieve enumerates $7499$ primes in $S_*$, with largest
prime $76163$, and their product has bit length $109445$. Hence
$B_*<2^{109445}$. With $A_0=4608^{165888}$, we have
$6^{100}<A_0$ and $128<A_0$, whence $H<8A_0B_*^{24}$. Since $4608<2^{13}$,
\[
 \log_2H<3+165888\cdot13+24\cdot109445=4783227.
\]
This is a check of the number H, not a proof that H is a height constant.

For completeness, the scalar product used by the original polynomial is
\[
 r(\beta)=\beta(\beta-1)(\beta-2401)(32\beta+343)(2\beta-1).
\]
At a prime $p>76176$, if $v_p(\beta)<0$, its valuation is
$5v_p(\beta)$ and so detects odd contact with infinity. If $\beta$ is
$p$-integral and has odd contact with a finite member of $\mathcal C$,
only that factor has positive valuation; every other factor is a unit.
The extra root $1/2$ introduces only exceptional collision primes
(notably factors of $4801$ and $718$). Thus every prime of $M(\beta)$ has
odd valuation in $r(\beta)$. If a separate comparison proves
$M(\beta)^{48}\mid s\ne0$, the five-contact estimate just proved gives
$h(\beta)\le H\sqrt{|s|}$.

Combining the non-CM and CM estimates with \eqref{eq:h5-R} and
\eqref{eq:h5-G} gives
\[
 h(s)\le 2(4608^{165888}+6^{100})(B_*M(s))^{24}+128
       \le H_*M(s)^{24},
\]
because $M(s)\ge1$. The numerical certificate above proves the
remaining assertion about $H_*$.
\end{proof}


